%! AP Statistics — Unit 3, Lesson 3.7 — Activity KEY
\documentclass[10pt]{article}
\usepackage{apstats-article}
\usepackage{apstats-key}

\begin{document}

\pageheader{Unit 3, Lesson 3.7}{Group Activity KEY: Scope of Inference Stations}
\namepartnerperiod

\vspace{0.1in}

\textbf{Study A:} Mindfulness app with volunteers.
\begin{practicebox}
\small
\begin{enumerate}[leftmargin=1.5em, itemsep=5pt]
  \item Random assignment: \ans{Yes.} \quad Random selection: \ans{No (volunteers).}
  \item Causation: \ans{Yes.} \quad Generalize to all college students: \ans{No.}
  \item \ans{Because random assignment was used, we can conclude that the mindfulness app caused the reduction in stress. However, because the participants were volunteers (not randomly selected), we cannot generalize this finding to all college students.}
  \item \ans{Randomly select college students from a larger population (e.g., from a college's registration list) rather than recruiting volunteers.}
\end{enumerate}
\end{practicebox}

\vspace{0.08in}

\textbf{Study B:} Social media observational study.
\begin{practicebox}
\small
\begin{enumerate}[leftmargin=1.5em, itemsep=5pt]
  \item Random assignment: \ans{No.} \quad Random selection: \ans{Yes.}
  \item Causation: \ans{No.} Explain: \ans{This is an observational study — no treatments were imposed, so we cannot rule out confounding variables.}
  \item \ans{Yes — the random sample allows generalization to the national adult population.}
  \item \ans{Depression, social isolation, age (any plausible confounder).}
\end{enumerate}
\end{practicebox}

\vspace{0.08in}

\textbf{Study C:} Drug trial with both types of randomization.
\begin{practicebox}
\small
\begin{enumerate}[leftmargin=1.5em, itemsep=5pt]
  \item Random assignment: \ans{Yes.} \quad Random selection: \ans{Yes.}
  \item Causation: \ans{Yes.} \quad Generalize: \ans{Yes.}
  \item \ans{To all patients at that hospital (the sampling frame), and potentially to patients similar to those at the hospital.}
  \item \ans{Because patients were randomly selected and randomly assigned, we can conclude that the new drug caused significantly less pain AND generalize this finding to the population of patients at the hospital.}
\end{enumerate}
\end{practicebox}

\vspace{0.08in}

\begin{practicebox}
\small
\textbf{Group Synthesis:}
\ans{The conclusion is partially wrong. Random assignment supports causation, but without random selection, we cannot generalize to all adults.}

Corrected conclusion: \ans{Because random assignment was used, we can conclude that the treatment caused the improvement. However, because participants were not randomly selected from the broader adult population, we cannot generalize these findings to all adults — only to individuals similar to those in the study.}
\end{practicebox}

\end{document}
